\documentclass[11pt,a4paper]{article}
\usepackage[utf8]{inputenc}
\usepackage[T1]{fontenc}
\usepackage[french]{babel}
\usepackage[margin=2.0cm]{geometry}
\usepackage{enumitem}
\usepackage{xcolor}
\usepackage{tcolorbox}
\usepackage{booktabs}
\usepackage{titlesec}
\usepackage{fancyhdr}
\usepackage{hyperref}
\usepackage{multicol}
\usepackage{tabularx}

% --- Couleurs BNP-inspired ---
\definecolor{bnpgreen}{HTML}{00915A}
\definecolor{bnpdark}{HTML}{1A1A2E}
\definecolor{accent}{HTML}{2563EB}
\definecolor{warmgray}{HTML}{4B5563}
\definecolor{lightbg}{HTML}{F0FDF4}
\definecolor{warnbg}{HTML}{FEF3C7}
\definecolor{infobg}{HTML}{EFF6FF}
\definecolor{darkbg}{HTML}{F3F4F6}
\definecolor{successbg}{HTML}{ECFDF5}
\definecolor{dangerbg}{HTML}{FEF2F2}
\definecolor{dangerframe}{HTML}{DC2626}

\hypersetup{colorlinks=true, linkcolor=accent, urlcolor=accent}

\tcbuselibrary{breakable}

% --- Boxes ---
\newtcolorbox{pitchbox}{
  colback=lightbg, colframe=bnpgreen, breakable,
  fonttitle=\bfseries\large, title={\textcolor{bnpgreen}{Pitch principal} \hfill \textit{$\sim$90 secondes}},
  arc=3pt, boxrule=1pt, left=8pt, right=8pt, top=6pt, bottom=6pt
}

\newtcolorbox{modulebox}[1]{
  colback=infobg, colframe=accent, breakable,
  fonttitle=\bfseries, title={#1},
  arc=2pt, boxrule=0.8pt, left=6pt, right=6pt, top=4pt, bottom=4pt
}

\newtcolorbox{warnbox}[1]{
  colback=warnbg, colframe={orange!80!black}, breakable,
  fonttitle=\bfseries, title={#1},
  arc=2pt, boxrule=0.8pt, left=6pt, right=6pt, top=4pt, bottom=4pt
}

\newtcolorbox{verbatimbox}{
  colback=white, colframe=warmgray, breakable,
  arc=2pt, boxrule=0.5pt, left=8pt, right=8pt, top=6pt, bottom=6pt
}

\newtcolorbox{goodbox}[1]{
  colback=successbg, colframe=bnpgreen, breakable,
  fonttitle=\bfseries, title={#1},
  arc=2pt, boxrule=0.8pt, left=6pt, right=6pt, top=4pt, bottom=4pt
}

\newtcolorbox{dangerbox}[1]{
  colback=dangerbg, colframe=dangerframe, breakable,
  fonttitle=\bfseries, title={#1},
  arc=2pt, boxrule=0.8pt, left=6pt, right=6pt, top=4pt, bottom=4pt
}

\newtcolorbox{graybox}[1]{
  colback=darkbg, colframe=warmgray, breakable,
  fonttitle=\bfseries, title={#1},
  arc=2pt, boxrule=0.8pt, left=6pt, right=6pt, top=4pt, bottom=4pt
}

% --- Headers ---
\titleformat{\section}{\Large\bfseries\color{bnpdark}}{}{0pt}{}[\color{bnpgreen}\rule{\textwidth}{0.8pt}]
\titleformat{\subsection}{\large\bfseries\color{accent}}{}{0pt}{}
\titleformat{\subsubsection}{\normalsize\bfseries\color{warmgray}}{}{0pt}{}

\pagestyle{fancy}
\fancyhf{}
\fancyhead[L]{\small\textcolor{warmgray}{Pr\'eparation entretien --- BNP Paribas CIB}}
\fancyhead[R]{\small\textcolor{warmgray}{CSPM --- Stage Analyste Cr\'edit}}
\fancyfoot[C]{\small\textcolor{warmgray}{\thepage}}
\renewcommand{\headrulewidth}{0.4pt}

\begin{document}

% ============================================================
%                       PAGE DE TITRE
% ============================================================
\begin{center}
{\LARGE\bfseries\textcolor{bnpdark}{Pr\'eparation Entretien}}\\[6pt]
{\Large\textcolor{bnpgreen}{BNP Paribas CIB --- Stage Assistant Analyste Cr\'edit}}\\[4pt]
{\large\textcolor{warmgray}{\'Equipe CSPM $\cdot$ Services MFS $\cdot$ Front Office $\cdot$ Pantin}}\\[12pt]
\textcolor{warmgray}{\rule{0.5\textwidth}{0.5pt}}\\[8pt]
{\small\textcolor{warmgray}{R\'ef. IS\_CIB2SIDF26\_0018 $\cdot$ 6 mois \`a partir de mars 2026}}
\end{center}

\vspace{4pt}

\tableofcontents

\newpage

% %%%%%%%%%%%%%%%%%%%%%%%%%%%%%%%%%%%%%%%%%%%%%%%%%%%%%%%%%%%%
%
%         PARTIE I --- PR\'ESENTATION DU PROJET
%
% %%%%%%%%%%%%%%%%%%%%%%%%%%%%%%%%%%%%%%%%%%%%%%%%%%%%%%%%%%%%

\part*{Partie I --- Pr\'esentation du projet}
\addcontentsline{toc}{part}{Partie I --- Pr\'esentation du projet}

% ============================================================
\section{Pitch principal --- Pr\'esentation du projet}
% ============================================================

\begin{pitchbox}
\og Mon m\'emoire de M1 \`a Paris-Saclay porte sur le \textbf{monitoring de fonds
et la gestion du collat\'eral} : j'ai construit un syst\`eme d'aide \`a la d\'ecision
pour des hedge funds, avec calcul automatis\'e de NAV et de LTV, suivi des margin calls,
valorisation et haircuts du collat\'eral, et stress tests de liquidit\'e sur portefeuilles.

En parall\`ele, j'ai d\'evelopp\'e un \textbf{cockpit de risque cr\'edit IFRS\,9} complet
--- calcul d'ECL, mod\`eles de probabilit\'e de d\'efaut, stress tests macro\'economiques,
RAROC et allocation de capital --- qui m'a donn\'e une \textbf{vision bout-en-bout} de la
cha\^ine de risque, du client individuel jusqu'au pilotage CRO.

Ces deux projets convergent vers le poste CSPM : le premier m'a donn\'e les
\textbf{outils Front Office} --- LTV, collat\'eral, margin calls --- et le second la
\textbf{profondeur analytique c\^ot\'e Risques}. C'est exactement l'intersection
du poste, \`a la fronti\`ere entre le Front Office et le d\'epartement des risques. \fg{}
\end{pitchbox}

\begin{warnbox}{Timing \& d\'ebit}
\begin{itemize}[nosep, leftmargin=12pt]
\item \textbf{90 secondes max} --- ne pas d\'epasser 2 minutes, l'interviewer doit pouvoir rebondir.
\item D\'ebit pos\'e, pauses naturelles apr\`es chaque bloc (m\'emoire / cockpit / convergence).
\item \textbf{Ne pas r\'eciter} --- adapter le wording selon la dynamique.
\item Si on te demande \og Pr\'esentez-vous \fg{} au sens large, ins\'erer le pitch \textit{apr\`es}
      2--3 phrases de parcours (Saclay, M1 finance/\'eco, quanti).
\end{itemize}
\end{warnbox}

\vspace{10pt}

% ============================================================
\section{Modules th\'ematiques --- Selon la conversation}
% ============================================================

\textcolor{warmgray}{\textit{Chaque module est autonome. Utilise celui qui correspond
\`a la direction prise par l'interviewer. Ne les d\'eroule \textbf{pas} tous --- choisis
1--2 modules pertinents selon les questions.}}

\vspace{8pt}

% --- Module A ---
\begin{modulebox}{Module A --- Si on te parle de \textbf{LTV / Collat\'eral / Margin Calls}}

\textit{C'est le c\oe{}ur du poste. C'est l\`a que ton m\'emoire M1 brille.}

\begin{verbatimbox}
\og Dans mon m\'emoire, j'ai impl\'ement\'e un monitoring quotidien de LTV avec
identification automatique des seuils critiques. Concr\`etement, le syst\`eme
calcule la NAV des fonds, en d\'eduit le LTV de chaque ligne de financement,
et d\'eclenche des alertes quand on approche des covenants.

Pour le collat\'eral, j'ai travaill\'e sur la valorisation des actifs mis en
garantie, le calcul des haircuts selon la nature de l'actif, et l'analyse
d'\'eligibilit\'e. J'ai aussi automatis\'e le suivi des margin calls ---
identification du d\'eficit, calcul du montant \`a appeler, et suivi de la
r\'egularisation. \fg{}
\end{verbatimbox}

\textbf{Mots-cl\'es \`a placer} : LTV, NAV, haircut, \'eligibilit\'e, margin call,
covenant, seuil critique, monitoring quotidien.

\textbf{Question de rebond possible} : \og Comment g\'eriez-vous un client dont le
LTV d\'epasse le seuil ? \fg{} $\rightarrow$ Escalade, appel de marge,
restructuration du collat\'eral, coordination avec le risk management.
\end{modulebox}

\vspace{6pt}

% --- Module B ---
\begin{modulebox}{Module B --- Si on te parle de \textbf{Risque de cr\'edit / Analyse financi\`ere}}

\textit{C'est l\`a que le projet IFRS\,9 te diff\'erencie des autres candidats.}

\begin{verbatimbox}
\og Mon cockpit IFRS\,9 calcule l'ECL --- Expected Credit Loss --- en trois
\'etapes : mod\'elisation de la PD, estimation de la LGD cyclique, et staging
SICR multi-facteurs. J'ai impl\'ement\'e trois familles de mod\`eles de PD
--- r\'egression logistique avec scorecard, XGBoost, et deep learning
tabulaire --- pour comparer leurs performances et r\'epondre aux exigences
de model risk management.

Ce qui est directement utile pour CSPM, c'est que je comprends comment
le risque assum\'e par la banque sur chaque exposition est couvert par la
garantie d\'etenue. Quand l'\'equipe CSPM \'echange avec le d\'epartement
des risques, je parle le m\^eme langage. \fg{}
\end{verbatimbox}

\textbf{Mots-cl\'es \`a placer} : ECL, PD, LGD, IFRS\,9, staging, SICR,
risk-weighted assets, model risk management.

\textbf{Attention} : ne pas tomber dans le technique pur. L'interviewer est
Front Office, pas quant. Rester sur la \textbf{logique m\'etier}.
\end{modulebox}

\vspace{6pt}

% --- Module C ---
\begin{modulebox}{Module C --- Si on te parle de \textbf{Stress Tests / Sc\'enarios de crise}}

\textit{Les deux projets se compl\`etent ici.}

\begin{verbatimbox}
\og J'ai travaill\'e sur les stress tests dans les deux projets. C\^ot\'e
monitoring de fonds, j'ai mod\'elis\'e des sc\'enarios de crise de liquidit\'e
et leur impact sur les covenants --- typiquement, que se passe-t-il si la
NAV du fonds chute de 20\,\% en un mois.

C\^ot\'e IFRS\,9, j'ai impl\'ement\'e des reverse stress tests qui
identifient les combinaisons macro\'economiques --- ch\^omage, taux,
immobilier --- qui am\`enent le portefeuille \`a un point de rupture.
C'est pertinent pour CSPM parce que l'\'equipe doit anticiper les
situations o\`u la capacit\'e de financement d'un client se d\'egrade
avant que \c{c}a ne devienne critique. \fg{}
\end{verbatimbox}

\textbf{Mots-cl\'es \`a placer} : sc\'enario adverse, reverse stress test,
point de rupture, anticipation, capacit\'e de financement.
\end{modulebox}

\vspace{6pt}

% --- Module D ---
\begin{modulebox}{Module D --- Si on te parle de \textbf{Python / Automatisation / Outils}}

\begin{verbatimbox}
\og J'ai automatis\'e l'ensemble de la cha\^ine avec Python : extraction
de donn\'ees march\'e, calculs de NAV et LTV, g\'en\'eration de reportings
Front Office. Mon cockpit IFRS\,9 recalcule un portefeuille complet de
30\,000 lignes en environ une seconde, avec export Excel multi-onglets.

J'utilise aussi VBA pour les reportings Excel classiques. Et de mani\`ere
plus large, j'ai d\'evelopp\'e une vraie m\'ethodologie de travail avec
l'IA --- pas juste des prompts basiques, mais du pilotage structur\'e
pour acc\'el\'erer l'impl\'ementation tout en gardant le contr\^ole sur
la qualit\'e et la rigueur financi\`ere. \fg{}
\end{verbatimbox}

\textbf{Mots-cl\'es \`a placer} : Python, VBA, automatisation, reporting,
Excel, gain de temps, scalabilit\'e.

\textbf{Transition naturelle vers le Module G} (IA) si l'interviewer rebondit.
\end{modulebox}

\vspace{6pt}

% --- Module E ---
\begin{modulebox}{Module E --- Si on te demande \textbf{\og Pourquoi BNP / Pourquoi CSPM ? \fg{}}}

\begin{verbatimbox}
\og Trois raisons. D'abord, le poste est \`a l'intersection exacte de mes
deux projets --- Front Office avec une dimension risque forte. Ensuite,
l'\'equipe CSPM travaille sur plusieurs offres de financement --- fonds de
couverture, ABC, Capdoc, FOF --- ce qui offre une exposition tr\`es large
aux produits de financement. Enfin, la relation directe avec les clients
et la coordination avec New York, Paris et Londres me motivent parce que
je veux comprendre le financement de bout en bout, pas seulement le
mod\`ele derri\`ere. \fg{}
\end{verbatimbox}

\textbf{Variante si on insiste} : mentionner le Graduate Programme et le
VIE comme perspectives d'\'evolution \`a moyen terme.
\end{modulebox}

\vspace{6pt}

% --- Module F ---
\begin{modulebox}{Module F --- Si on te demande tes \textbf{qualit\'es / d\'efauts / soft skills}}

\begin{verbatimbox}
\og Je suis rigoureux --- mon cockpit IFRS\,9 a 600 tests automatis\'es
et a pass\'e trois audits successifs. Je suis aussi autonome --- les deux
projets sont des initiatives personnelles, pas des travaux dirig\'es.
Et j'aime aller au fond des sujets : quand j'ai identifi\'e des
approximations math\'ematiques dans mon moteur de calcul, j'ai men\'e
un audit de 24 points et tout corrig\'e plut\^ot que de laisser passer.

Mon point d'am\'elioration, c'est que j'ai tendance \`a vouloir trop
optimiser avant de livrer. J'apprends \`a privil\'egier le pragmatisme
--- livrer une version fonctionnelle puis it\'erer. \fg{}
\end{verbatimbox}
\end{modulebox}

\vspace{6pt}

% --- Module G --- IA
\begin{modulebox}{Module G --- Si on te parle d'\textbf{IA} ou si on te demande \textbf{\og Comment tu as fait le projet ? \fg{}}}

\textit{Ce module est important. Le cockpit fait 15\,000 lignes --- un M1 ne code pas
\c{c}a seul en deux mois. Il faut \^etre honn\^ete et transformer \c{c}a en force.}

\begin{verbatimbox}
\og Le cockpit fait 15\,000 lignes de code, 600 tests, et a pass\'e trois
audits de rigueur math\'ematique. En M1, en deux mois, c'est \'evidemment
pas du code \'ecrit ligne par ligne \`a la main. J'ai utilis\'e l'IA de
mani\`ere intensive --- pas juste des prompts basiques, mais une vraie
m\'ethodologie de pilotage.

Concr\`etement : je con\c{c}ois l'architecture en amont, je r\'edige des
sp\'ecifications fonctionnelles d\'etaill\'ees, je pilote des agents IA
pour l'impl\'ementation, et surtout je fais la revue, la validation et la
correction de chaque module.

Ma valeur ajout\'ee c'est pas d'\'ecrire du Python --- c'est de
\textbf{savoir qu'un ECL se calcule avec un taux de hasard et pas une
approximation de Bernoulli}, que la LGD doit \^etre cyclique, que le RAROC
int\`egre un CIR et un taux d'imposition. L'IA impl\'emente vite, mais
elle ne sait pas si c'est financièrement correct. Moi si.

Et concr\`etement pour une \'equipe comme CSPM, \c{c}a veut dire que l\`a o\`u
quelqu'un met trois mois \`a automatiser un reporting ou un suivi de LTV,
je peux le faire en une semaine. C'est un vrai levier de productivit\'e. \fg{}
\end{verbatimbox}

\begin{goodbox}{Ce que ce module d\'emontre}
\begin{enumerate}[nosep, leftmargin=16pt]
\item \textbf{Honn\^etet\'e} --- tu ne pr\'etends pas avoir tout cod\'e \`a la main
\item \textbf{Expertise m\'etier} --- tu sais ce qui est financièrement correct
\item \textbf{M\'ethodologie} --- ce n'est pas \og copier-coller de ChatGPT \fg{}
\item \textbf{Valeur concr\`ete pour l'\'equipe} --- gain de productivit\'e mesurable
\end{enumerate}
\end{goodbox}

\begin{dangerbox}{Les lignes rouges}
\begin{itemize}[nosep, leftmargin=12pt]
\item \textbf{Ne dis jamais} \og l'IA a cod\'e mon moteur d'ECL \fg{}
      --- dis \og j'ai con\c{c}u et valid\'e un moteur d'ECL, en utilisant l'IA comme acc\'el\'erateur \fg{}
\item \textbf{Ne d\'etaille pas} l'outillage (Claude Code, sous-agents, orchestration)
      --- l'interviewer s'en fiche, il veut savoir ce que \c{c}a produit
\item \textbf{Ne t'excuse pas} --- pr\'esente-le comme une comp\'etence, pas une b\'equille
\end{itemize}
\end{dangerbox}

\textbf{M\'etaphore utile si besoin} : \og C'est comme un architecte qui ne pose pas
les briques lui-m\^eme mais qui con\c{c}oit le b\^atiment et v\'erifie qu'il tient debout. \fg{}

\textbf{Rebond offensif} : \og Si l'\'equipe CSPM a besoin d'automatiser un nouveau
suivi de LTV ou de cr\'eer un reporting pour un type de fonds, je sais comment
sp\'ecifier le besoin, produire rapidement avec l'IA, et surtout valider que le
r\'esultat est financièrement correct. C'est pas juste taper un prompt et copier-coller. \fg{}

\end{modulebox}

\vspace{10pt}

% ============================================================
\section{Correspondance poste $\leftrightarrow$ comp\'etences}
% ============================================================

\begin{center}
\small
\begin{tabular}{@{}p{5.2cm}p{4.0cm}p{4.0cm}@{}}
\toprule
\textbf{Mission CSPM} & \textbf{M\'emoire M1} & \textbf{Cockpit IFRS\,9} \\
\midrule
V\'erifier niveaux de LTV & Calcul NAV \& LTV, seuils & --- \\[3pt]
Monitoring collat\'eral & Haircuts, \'eligibilit\'e & LGD cyclique, collateral \\[3pt]
Analyser margin calls & Suivi margin calls auto. & --- \\[3pt]
Optimiser financement & --- & RAROC, allocation capital \\[3pt]
Relation avec Risques & --- & ECL, staging, stress tests \\[3pt]
Onboarding clients & --- & --- \textit{(motivation)} \\[3pt]
Stress tests & Tests liquidit\'e, covenants & RST, sc\'enarios macro \\[3pt]
Excel, AI, automatisation & Python/VBA, reportings & Pipeline 1s, IA pilot\'ee \\
\bottomrule
\end{tabular}
\end{center}


% %%%%%%%%%%%%%%%%%%%%%%%%%%%%%%%%%%%%%%%%%%%%%%%%%%%%%%%%%%%%
%
%         PARTIE II --- GUIDE G\'EN\'ERAL D'ENTRETIEN
%
% %%%%%%%%%%%%%%%%%%%%%%%%%%%%%%%%%%%%%%%%%%%%%%%%%%%%%%%%%%%%

\newpage
\part*{Partie II --- Guide g\'en\'eral d'entretien}
\addcontentsline{toc}{part}{Partie II --- Guide g\'en\'eral d'entretien}

% ============================================================
\section{BNP Paribas --- Histoire \& actualit\'e CIB}
% ============================================================

\subsection{Dates cl\'es --- De la cr\'eation \`a aujourd'hui}

\begin{graybox}{Chronologie}
\begin{description}[style=nextline, leftmargin=0.8cm, nosep, itemsep=5pt]

\item[\textbf{1848}] Cr\'eation du Comptoir National d'Escompte de Paris (CNEP) ---
l'une des plus anciennes institutions bancaires fran\c{c}aises.

\item[\textbf{1872}] Fondation de la Banque de Paris et des Pays-Bas (\textbf{Paribas}),
banque d'investissement \`a vocation internationale.

\item[\textbf{1966}] Naissance de la \textbf{BNP} (Banque Nationale de Paris) par fusion
du CNEP et de la BNCI (Banque Nationale pour le Commerce et l'Industrie).

\item[\textbf{1993}] Privatisation de la BNP par l'\'Etat fran\c{c}ais.

\item[\textbf{1999--2000}] \textbf{Naissance de BNP Paribas.}
F\'evrier 1999 : la Soci\'et\'e G\'en\'erale annonce une fusion avec Paribas.
Mars 1999 : BNP lance une contre-offre hostile sur \textit{les deux} --- une premi\`ere
dans la banque fran\c{c}aise. BNP obtient 65,2\,\% de Paribas mais seulement 31,8\,\%
de la SocGen. R\'esultat : BNP + Paribas fusionnent le \textbf{23 mai 2000}.
La nouvelle entit\'e devient la premi\`ere banque fran\c{c}aise et parmi le top\,5 europ\'een.

\item[\textbf{2009}] \textbf{Acquisition de Fortis Bank} (Belgique \& Luxembourg) pour
14,5\,Md\texteuro{}. Op\'eration transform\-ationnelle men\'ee en pleine crise financi\`ere
apr\`es l'effondrement de Fortis. BNP Paribas devient la \textbf{premi\`ere banque de Belgique}.

\item[\textbf{2019}] \textbf{Rachat du Prime Brokerage de Deutsche Bank} + Electronic Equities.
$\sim$900 collaborateurs transf\'er\'es dans le monde. Donne \`a BNPP une plateforme
d'ex\'ecution et de risk management de premier plan, pris\'ee par les fonds quantitatifs.
Int\'egration achev\'ee fin 2021. \textit{C'est le moment fondateur du MFS actuel.}

\item[\textbf{2022}] Vente de Bank of the West (banque de d\'etail US) \`a BMO pour
$\sim$16,3\,Md\$. BNP se recentre sur le CIB et la banque de d\'etail europ\'eenne.
Lancement du plan strat\'egique \textbf{GTS\,2025} (Growth, Technology, Sustainability).

\item[\textbf{2023}] \textbf{Effondrement de Credit Suisse.} BNP recrute des \'equipes
seniors de CS en prime brokerage et capte le flux client via un accord de r\'ef\'erencement.
Opportunit\'e g\'en\'erationnelle pour gagner des parts de march\'e.

\item[\textbf{2024--2025}] Acquisition d'\textbf{AXA Investment Managers} pour
5,1\,Md\texteuro{} (finalis\'ee juillet 2025). Cr\'ee un leader europ\'een de la gestion
d'actifs avec 1\,500\,Md\texteuro{} d'encours.

\end{description}
\end{graybox}

\subsection{Chiffres cl\'es du Groupe (au 31 d\'ecembre 2025)}

\begin{center}
\small
\begin{tabular}{@{}ll@{}}
\toprule
\textbf{Indicateur} & \textbf{Valeur} \\
\midrule
Total des actifs & $\sim$2\,800\,Md\texteuro{} (2\textsuperscript{e} banque europ\'eenne, 8\textsuperscript{e} mondiale) \\
Capitalisation boursi\`ere & $\sim$105\,Md\texteuro{} (jan.\,2026) \\
Revenus Groupe (2025) & 51,2\,Md\texteuro{} (+4,9\,\%) \\
R\'esultat net (2025) & \textbf{12,2\,Md\texteuro{}} --- record historique (+4,6\,\%) \\
Ratio CET1 & 12,6\,\% \\
ROTE & 11,6\,\% \\
Collaborateurs & $\sim$190\,000 dans 64 pays \\
Dividende 2025 & 5,16\,\texteuro{}/action (+7,7\,\%), payout 60\,\% \\
\bottomrule
\end{tabular}
\end{center}

\subsection{CIB --- R\'esultats 2025 (publi\'es le 5 f\'evrier 2026)}

\textit{Le CIB a r\'ealis\'e une \textbf{ann\'ee record} avec +4,5\,\% de revenus
(+7,7\,\% \`a taux de change constants).}

\begin{center}
\small
\begin{tabular}{@{}lrr@{}}
\toprule
\textbf{M\'etier CIB} & \textbf{Revenus 2025} & \textbf{$\Delta$ YoY} \\
\midrule
Global Banking & 6\,244\,M\texteuro{} & $-$0,5\,\% \\
Global Markets --- FICC & 5\,522\,M\texteuro{} & +8,3\,\% \\
\textbf{Global Markets --- Equity \& Prime Services} & \textbf{4\,046\,M\texteuro{}} & \textbf{+10,2\,\%} \\
Securities Services & 3\,185\,M\texteuro{} & +8,1\,\% \\
\midrule
\textbf{CIB Total} & \textbf{$\sim$19,0\,Md\texteuro{}} & \textbf{+4,5\,\%} \\
\bottomrule
\end{tabular}
\end{center}

\begin{warnbox}{\`A retenir pour l'entretien}
\textbf{Equity \& Prime Services} (l\`a o\`u se situe CSPM) est le m\'etier CIB qui
cro\^it le plus vite : +10,2\,\% en 2025, +42,1\,\% au T1\,2025.
Les encours prime brokerage ont \textbf{quadrupl\'e en 6 ans}.
CAGR revenus >35\,\% depuis l'acquisition de Deutsche Bank en 2019.
\end{warnbox}

\subsection{Dirigeants cl\'es}

\begin{center}
\small
\begin{tabular}{@{}p{5cm}p{4cm}p{4.5cm}@{}}
\toprule
\textbf{Fonction} & \textbf{Nom} & \textbf{Contexte} \\
\midrule
DG Groupe & Jean-Laurent Bonnaf\'e & DG depuis 2011, reconduit mai 2025 \\
DGA + Pr\'esident CIB & Yann G\'erardin & Nomm\'e sept.\,2025 \\
DGA + DG CIB & Olivier Osty & Ex-Head Global Markets \\
Head Prime Brokerage & Ashley Wilson & Ex-Deutsche Bank \\
\bottomrule
\end{tabular}
\end{center}

\subsection{Plan strat\'egique GTS\,2025 \& objectifs 2028}

\begin{graybox}{GTS\,2025 --- 3 piliers}
\begin{description}[style=nextline, leftmargin=0.8cm, nosep, itemsep=4pt]
\item[\textbf{Growth}] Consolider le leadership europ\'een, croissance organique rentable.
\item[\textbf{Technology}] IA, APIs, cloud comme fondations. \textbf{800+ applications IA en production}.
Plateforme \og LLM as a Service \fg{} interne (partenariat strat\'egique avec Mistral AI).
\item[\textbf{Sustainability}] 252\,Md\texteuro{} de financements durables accord\'es en 2025
(objectif d\'epass\'e). Cadre Green Bond mis \`a jour en mai 2025.
\end{description}
\end{graybox}

\begin{goodbox}{Objectifs 2028 (relev\'es en f\'evrier 2026)}
\begin{itemize}[nosep, leftmargin=12pt]
\item ROTE cible : \textbf{>13\,\%} (relev\'e)
\item Croissance r\'esultat net : \textbf{CAGR >10\,\%} (2025--2028)
\item Ratio CET1 cible : \textbf{13\,\%} fin 2027
\item Coefficient d'exploitation : \textbf{<56\,\%} d'ici 2028
\item \'Economies op\'erationnelles : 600\,M\texteuro{} en 2025 + 600\,M\texteuro{} en 2026
\end{itemize}
\end{goodbox}

\subsection{Prime Brokerage \& Fund Financing --- L'\'ecosyst\`eme CSPM}

\begin{graybox}{\'Evolution du Prime Brokerage chez BNP}
\begin{enumerate}[nosep, leftmargin=16pt]
\item \textbf{Avant 2019} : prime brokerage de taille moyenne, principalement europ\'een,
forte franchise FICC mais equity prime limit\'e.
\item \textbf{2019 --- Acquisition Deutsche Bank} : plateforme d'ex\'ecution \'electronique
reconnue, $\sim$900 collaborateurs, relations \'etablies avec les grands hedge funds
multi-strat\'egies et quantitatifs. \textit{Moment fondateur.}
\item \textbf{2021--2022} : Int\'egration compl\`ete. Plateforme prime mondiale unifi\'ee.
\item \textbf{2023} : Sortie de Credit Suisse --- BNP recrute des \'equipes CS seniors et
capture le flux client. Opportunit\'e unique.
\item \textbf{2024--2026} : +600\,bp de part de portefeuille chez les 30 premiers clients.
Expansion Asie, Moyen-Orient, US.
\end{enumerate}
\end{graybox}

\begin{graybox}{Offre de financement de fonds --- ce que fait MFS}
\begin{itemize}[nosep, leftmargin=12pt]
\item \textbf{Equity bridge facilities} (lignes de souscription / capital call facilities)
      adoss\'ees aux engagements non appel\'es des LPs
\item \textbf{Margin financing} dans 31 march\'es et 13 devises
\item \textbf{Securities lending} avec inventaire profond (y compris Chine, march\'es du Golfe)
\item \textbf{Cross-margining} actions, cr\'edit, futures et options sur une plateforme unique
\item \textbf{Capital introduction} --- \'equipes d\'edi\'ees dans toutes les r\'egions
\end{itemize}
BNP est l'une des \textbf{4 seules banques au monde} offrant la suite compl\`ete
de services aux investisseurs (prime brokerage + clearing + d\'eriv\'es list\'es +
repo/financing).
\end{graybox}

\subsection{Paysage concurrentiel --- O\`u se situe BNP}

\begin{center}
\small
\begin{tabular}{@{}clp{6.5cm}@{}}
\toprule
\textbf{Tier} & \textbf{Banque} & \textbf{Position} \\
\midrule
\multirow{3}{*}{\textbf{Tier 1}} & Goldman Sachs & \#1 mondial, $\sim$63\,\% de p\'en\'etration chez les HF >1\,Md\$ \\
 & Morgan Stanley & \#2, 226 nouveaux fonds au T1\,2025 \\
 & JP Morgan & \#3, 210 nouveaux fonds au T1\,2025 \\
\midrule
\multirow{3}{*}{\textbf{Tier 2}} & \textbf{BNP Paribas} & \textbf{\#6--7 mondial, \#1--2 en Europe},
  challenger \`a la croissance la plus rapide (CAGR >35\,\%) \\
 & Barclays & Fort en equities \& FICC prime europ\'een \\
 & UBS (ex-CS) & En phase de reconsolidation post-int\'egration CS \\
\bottomrule
\end{tabular}
\end{center}

\textit{Ambition d\'eclar\'ee de BNP : devenir \og incontestablement le premier en dehors
des US \fg{}.}

\subsection{R\'eglementation --- B\^ale\,IV / CRR3}

\begin{graybox}{Ce qu'il faut savoir}
\begin{itemize}[nosep, leftmargin=12pt]
\item \textbf{CRR3 en vigueur depuis le 1\textsuperscript{er} janvier 2025} dans l'UE,
      CRD6 depuis le 11 janvier 2026
\item \textbf{Output floor} : 50\,\% en 2025, mont\'ee progressive \`a 72,5\,\% d'ici 2030
\item \textbf{Impact sur BNP} : $\sim$20\,bp de hausse des exigences de capital vs. pairs.
      Le ratio CET1 peut baisser \`a mesure que les exigences montent.
\item \textbf{FRTB} (risque de march\'e) : certains \'el\'ements report\'es au 1\textsuperscript{er}
      janvier 2026
\item BNP cible un CET1 de 13\,\% fin 2027, offrant un coussin substantiel
\end{itemize}
\textbf{Lien avec CSPM} : CRR3 augmente le co\^ut en capital des expositions.
Un bon monitoring du collat\'eral (LTV bas, haircuts conservateurs) \textbf{r\'eduit le RW}
et donc le capital immobilis\'e. Le travail de CSPM a un impact direct sur l'efficacit\'e
du bilan.
\end{graybox}

\subsection{IA chez BNP Paribas --- \'El\'ements \`a placer}

\begin{graybox}{Les initiatives IA du Groupe}
\begin{itemize}[nosep, leftmargin=12pt]
\item \textbf{800+ applications IA en production} dans tous les m\'etiers
\item Plateforme \textbf{\og LLM as a Service \fg{}} : acc\`es s\'ecuris\'e aux mod\`eles
      open source, Mistral AI et mod\`eles internes via une interface standardis\'ee
\item \textbf{NOA} : assistant IA d\'eploy\'e sur la plateforme Centric pour les clients
      corporate du CIB
\item \textbf{CapLink Private} : portail automatis\'e par IA/robotique pour le traitement
      et la validation des appels de capital chez les clients fonds
\item \textbf{Outils ESG} : NLP/LLM pour estimer les \'emissions GES (Scope 1 \& 2)
      des entreprises non-reportantes
\item \textbf{Intelligent Document Processing} : automatisation des workflows documentaires
\end{itemize}
\end{graybox}

\subsection{Distinctions r\'ecentes}

\begin{itemize}[nosep, leftmargin=12pt]
\item \textbf{IFR Awards 2025} : 5 prix dont Sustainable Finance House of the Year,
      Euro Bond House of the Year, EMEA CLO House (2\textsuperscript{e} ann\'ee cons\'ecutive)
\item \textbf{Euromoney 2025} : 14 prix dont Europe's Best Bank for Securities Services,
      Europe's Best Investment Bank for Financing
\item \textbf{Risk.net} : \textbf{Prime Broker of the Year} --- BNP Paribas
\item \textbf{\#1 Western Europe Bookrunner 2025} : 48,8\,Md\$ de volume, 5,33\,\% de
      part de march\'e --- devant JP Morgan
\item \textbf{\#1 EMEA Loan Bookrunner 2024}
\end{itemize}

\vspace{6pt}

\begin{goodbox}{Synth\`ese --- Les 5 faits \`a conna\^itre sur BNP CIB}
\begin{enumerate}[nosep, leftmargin=16pt]
\item \textbf{Premi\`ere banque de la zone euro}, 190\,000 collaborateurs, 64 pays
\item \textbf{Equity \& Prime Services} = m\'etier CIB qui cro\^it le plus vite (+10,2\,\% en 2025),
      quadruplement des encours prime en 6 ans
\item \textbf{Acquisition Deutsche Bank 2019 + exit Credit Suisse 2023} = moment
      g\'en\'erationnel pour le fund financing
\item \textbf{800+ applications IA en production}, plateforme LLM interne, partenariat Mistral
\item \textbf{R\'esultat net record 2025} (12,2\,Md\texteuro{}), objectifs 2028 relev\'es
\end{enumerate}
\end{goodbox}

\vspace{10pt}

% ============================================================
\section{Comprendre le poste --- Ce que fait l'\'equipe CSPM}
% ============================================================

\begin{graybox}{L'\'equipe CSPM en bref}
\begin{itemize}[nosep, leftmargin=12pt]
\item \textbf{Rattachement} : Services Clients \& Gestion de Portefeuille, au sein du
      Front Office des Services MFS (Margin Financing Services)
\item \textbf{Activit\'e} : financement de fonds --- \`a l'origine fonds de fonds (FoF) avec
      faible LTV (Loan-to-Value), p\'erim\`etre \'elargi depuis
\item \textbf{Produits de financement} :
  \begin{itemize}[nosep, leftmargin=12pt]
    \item Fonds de couverture (hedge funds)
    \item Financement ABC (Asset-Based Credit --- adoss\'e aux actifs)
    \item Offre de discount engag\'e
    \item Capdoc (Capital Document --- facilit\'e de cr\'edit)
    \item Financement FOHF de tiers avec facilit\'es de levier
  \end{itemize}
\item \textbf{R\^ole cl\'e} : centraliser les demandes clients li\'ees au financement
      \textbf{et} fournir une assistance post-trade (options, change, FX, pr\^ets, remboursements)
\item \textbf{Position dans la banque} : Front Office mais relation forte avec le d\'epartement
      des risques --- le LTV est l'indicateur qui connecte les deux
\end{itemize}
\end{graybox}

\subsection{Les 6 missions du stagiaire --- d\'ecrypt\'ees}

\begin{enumerate}[leftmargin=16pt]
\item \textbf{V\'erifier le niveau de LTV des clients en financement au sein de notre desk}\\
\textcolor{warmgray}{$\rightarrow$ Tu re\c{c}ois des donn\'ees de NAV des fonds, tu calcules le LTV
(encours de pr\^et / valeur du collat\'eral), tu v\'erifies qu'on est sous le seuil contractuel.
C'est du monitoring quotidien, souvent sur Excel/Bloomberg.}

\item \textbf{Monitoring du collat\'eral et optimisation de la capacit\'e de financement}\\
\textcolor{warmgray}{$\rightarrow$ V\'erifier que les actifs en garantie sont toujours \'eligibles,
appliquer les bons haircuts, identifier des marges de man\oe{}uvre pour que le client puisse
emprunter davantage sans d\'epasser les seuils.}

\item \textbf{Analyser des demandes de marge et \'emettre des pre-approvals}\\
\textcolor{warmgray}{$\rightarrow$ Quand le LTV se d\'egrade, il faut appeler de la marge (margin call).
Tu analyses la situation, estimes le montant, et pr\'epares un avis pour le manager.}

\item \textbf{R\'epondre aux diverses demandes des clients}\\
\textcolor{warmgray}{$\rightarrow$ Questions sur les positions, les mouvements de tr\'esorerie,
les \'ech\'eanciers. Relation client directe.}

\item \textbf{Participation \`a l'onboarding de nouveaux clients}\\
\textcolor{warmgray}{$\rightarrow$ Analyse des portefeuilles du nouveau client avec la Structuration,
matrices d'\'eligibilit\'e, capacit\'e d'emprunt.}

\item \textbf{Demandes ad hoc de structuration et du management}\\
\textcolor{warmgray}{$\rightarrow$ Analyses ponctuelles, projets d'am\'elioration, reportings sp\'eciaux.
C'est l\`a que tes comp\'etences Python/IA sont les plus valoris\'ees.}
\end{enumerate}

\vspace{8pt}

% ============================================================
\section{Vocabulaire technique --- Ma\^itriser avant l'entretien}
% ============================================================

\textcolor{warmgray}{\textit{L'interviewer peut tester tes connaissances de base. Tu dois
d\'efinir chaque terme ci-dessous en 1--2 phrases, sans h\'esiter.}}

\subsection{Financement \& collat\'eral}

\begin{graybox}{D\'efinitions cl\'es}
\begin{description}[style=nextline, leftmargin=1cm, nosep, itemsep=6pt]
\item[\textbf{LTV (Loan-to-Value)}]
Ratio encours de pr\^et / valeur du collat\'eral.
Exemple : pr\^et de 70\,M\texteuro{} garanti par un portefeuille de 100\,M\texteuro{}
$\rightarrow$ LTV = 70\,\%.
Un LTV qui monte = risque qui augmente.

\item[\textbf{NAV (Net Asset Value)}]
Valeur liquidative nette d'un fonds = actifs $-$ passifs.
Publi\'ee quotidiennement ou mensuellement selon le fonds.
C'est la base du calcul de LTV.

\item[\textbf{Haircut}]
D\'ecote appliqu\'ee \`a la valeur d'un actif mis en garantie pour refl\'eter
sa volatilit\'e et son risque de liquidit\'e.
Exemple : actions liquides $\rightarrow$ haircut 10\,\% ; immobilier $\rightarrow$ haircut 40\,\%.
Collat\'eral apr\`es haircut = valeur march\'e $\times$ (1 $-$ haircut).

\item[\textbf{Margin call (appel de marge)}]
Demande adress\'ee au client de d\'eposer du collat\'eral suppl\'ementaire
(ou de rembourser une partie du pr\^et) quand le LTV d\'epasse le seuil contractuel.

\item[\textbf{Covenant}]
Clause contractuelle qui impose des limites au client (LTV max,
diversification minimale, etc.). Le non-respect peut d\'eclencher un
margin call ou un \'ev\'enement de d\'efaut.

\item[\textbf{\'Eligibilit\'e}]
Crit\`eres qui d\'eterminent si un actif peut servir de collat\'eral
(liquidit\'e, notation, devise, concentration\ldots).

\item[\textbf{ABC (Asset-Based Credit)}]
Financement adoss\'e directement aux actifs du fonds (pas \`a sa NAV globale).
Chaque actif a un taux d'avance sp\'ecifique selon son profil de risque.

\item[\textbf{Capdoc (Capital Call Document)}]
Facilit\'e de cr\'edit adoss\'ee aux engagements non appel\'es (\textit{unfunded commitments})
des investisseurs d'un fonds. Utilis\'e par les fonds de PE pour financer des acquisitions
avant d'appeler le capital des LPs.

\item[\textbf{FoF / FOHF (Fund of Hedge Funds)}]
Fonds qui investit dans un portefeuille de hedge funds.
Diversification mais double couche de frais. Historiquement le c\oe{}ur de CSPM.

\item[\textbf{Facility / ligne de cr\'edit}]
Accord de financement entre la banque et le client : montant maximal,
dur\'ee, collat\'eral requis, covenants. Le stagiaire monitore l'utilisation.
\end{description}
\end{graybox}

\subsection{Risque de cr\'edit (niveau Front Office)}

\begin{graybox}{Notions \`a conna\^itre --- pas besoin d'\^etre expert, mais de comprendre la logique}
\begin{description}[style=nextline, leftmargin=1cm, nosep, itemsep=6pt]
\item[\textbf{Exposition (EAD)}]
Montant que la banque perdrait si le client fait d\'efaut.
Pour CSPM = encours du pr\^et + engagements potentiels.

\item[\textbf{RWA (Risk-Weighted Assets)}]
Actifs pond\'er\'es par le risque --- base du calcul de capital r\'eglementaire.
Un LTV \'elev\'e $\rightarrow$ RW plus \'elev\'e $\rightarrow$ plus de capital immobilis\'e.

\item[\textbf{ECL (Expected Credit Loss)}]
Perte attendue sous IFRS\,9 = PD $\times$ LGD $\times$ EAD.
Tu n'as pas besoin de le calculer en stage, mais tu dois comprendre
que le d\'epartement des risques l'utilise pour \'evaluer les expositions
que tu monitores.

\item[\textbf{IFRS\,9 --- staging}]
Stage 1 (performant), Stage 2 (d\'egradation significative), Stage 3 (d\'efaut).
Un client dont le LTV se d\'egrade peut passer en Stage 2 $\rightarrow$ impact
provisionnel pour la banque.
\end{description}
\end{graybox}

\subsection{Environnement BNP CIB}

\begin{graybox}{Contexte organisationnel}
\begin{description}[style=nextline, leftmargin=1cm, nosep, itemsep=6pt]
\item[\textbf{CIB (Corporate \& Institutional Banking)}]
Banque de financement et d'investissement. C'est le p\^ole qui g\`ere
les grands clients institutionnels (fonds, corporates, souverains).

\item[\textbf{MFS (Margin Financing Services)}]
Sous-ensemble du Front Office qui fournit des solutions de financement
avec effet de levier aux fonds d'investissement.

\item[\textbf{Front Office vs. Middle/Back Office}]
Front Office = contact client, pricing, structuration, monitoring.
CSPM est FO mais avec une forte composante op\'erationnelle (monitoring quotidien).

\item[\textbf{Structuration}]
\'Equipe qui con\c{c}oit les produits de financement (termes, collat\'eral,
covenants). CSPM travaille en aval : une fois le deal structur\'e,
c'est CSPM qui monitore au quotidien.
\end{description}
\end{graybox}

\vspace{8pt}

% ============================================================
\section{Questions classiques --- R\'eponses pr\'epar\'ees}
% ============================================================

\subsection{\og Pr\'esentez-vous \fg{} ($\sim$2 minutes)}

\begin{verbatimbox}
\og Je suis en M1 \`a Paris-Saclay, sp\'ecialisation finance et \'economie
quantitative. Mon parcours est orient\'e vers l'analyse de risque et le
financement : mon m\'emoire porte sur le monitoring de fonds et la gestion
du collat\'eral, et j'ai d\'evelopp\'e en parall\`ele un cockpit de risque
cr\'edit IFRS\,9 complet.

\textit{[Ins\'erer le pitch de 90s]}

Je cherche un stage o\`u je peux appliquer ces comp\'etences dans un
environnement op\'erationnel, avec une relation client directe. Le poste
CSPM coche toutes les cases. \fg{}
\end{verbatimbox}

\subsection{\og Qu'est-ce que vous connaissez de BNP Paribas ? \fg{}}

\begin{verbatimbox}
\og BNP Paribas est la premi\`ere banque de la zone euro, avec un r\'esultat
net record de 12,2 milliards d'euros en 2025. C\^ot\'e CIB, le m\'etier
Equity \& Prime Services cro\^it de plus de 10\,\% par an --- c'est le
segment le plus dynamique, port\'e par l'acquisition du prime brokerage
de Deutsche Bank en 2019 et la captation de clients apr\`es la sortie
de Credit Suisse en 2023.

Ce qui m'int\'eresse sp\'ecifiquement, c'est que BNP est l'une des quatre
seules banques au monde \`a offrir la suite compl\`ete de services aux
investisseurs, et que l'\'equipe CSPM est au c\oe{}ur de cette dynamique
--- avec un p\'erim\`etre qui s'est \'elargi au-del\`a des fonds de fonds
initiaux pour couvrir le hedge fund financing, l'ABC et le Capdoc.

Et le plan GTS\,2025 montre une banque qui investit massivement dans
l'IA --- plus de 800 applications en production, un partenariat avec
Mistral --- ce qui est coh\'erent avec mon profil. \fg{}
\end{verbatimbox}

\subsection{\og Votre plus grand d\'efi / \'echec ? \fg{}}

\begin{verbatimbox}
\og Dans mon cockpit IFRS\,9, j'ai d\'ecouvert apr\`es des semaines de
d\'eveloppement que mon RAROC \'etait insensible aux stress
macro\'economiques --- il restait \`a 2,85\,\% quel que soit le sc\'enario.

Le probl\`eme \'etait subtil : les revenus d'int\'er\^ets \'etaient
calcul\'es \`a partir du spread stress\'e, ce qui compensait presque
exactement la hausse des pertes attendues. Il a fallu repenser le
mod\`ele en s\'eparant les spreads d'origination des conditions de
march\'e stress\'ees.

Ce que j'en retiens : quand un r\'esultat semble trop stable, il faut
creuser. La rigueur de validation est aussi importante que la mod\'elisation. \fg{}
\end{verbatimbox}

\subsection{\og O\`u vous voyez-vous dans 5 ans ? \fg{}}

\begin{verbatimbox}
\og \`A court terme, je veux consolider mes comp\'etences en financement
structur\'e --- comprendre les produits, la relation client, le monitoring
au quotidien. Si le stage se passe bien, le Graduate Programme ou un VIE
me permettrait de voir d'autres desks ou d'autres g\'eographies.

\`A moyen terme, je me vois \'evoluer vers un r\^ole d'analyste senior
en financement de fonds, \`a l'intersection entre la structuration et
la gestion du risque --- exactement l\`a o\`u se situe CSPM aujourd'hui. \fg{}
\end{verbatimbox}

\subsection{\og Pourquoi vous et pas un autre candidat ? \fg{}}

\begin{verbatimbox}
\og La plupart des candidats M1 connaissent la th\'eorie du financement.
Moi, j'ai impl\'ement\'e un syst\`eme de monitoring de LTV et de margin
calls de bout en bout. Et avec le cockpit IFRS\,9, je comprends aussi
le c\^ot\'e risque --- ECL, staging, capital --- ce qui me permet d'\^etre
op\'erationnel quand l'\'equipe \'echange avec le d\'epartement des risques.

En plus, je sais utiliser l'IA pour acc\'el\'erer le travail --- pas
juste des prompts, mais un pilotage structur\'e qui me rend concr\`etement
plus productif. Si l'\'equipe a besoin d'automatiser un nouveau reporting
ou un suivi, je peux livrer en jours, pas en semaines. \fg{}
\end{verbatimbox}

\subsection{\og Avez-vous d\'ej\`a travaill\'e en \'equipe ? \fg{}}

\begin{verbatimbox}
\og Mes projets sont individuels, mais ils m'ont appris des m\'ethodes
d'\'equipe. Mon cockpit a 600 tests automatis\'es et une architecture
modulaire --- c'est exactement ce qu'on fait quand on code \`a plusieurs
pour que chacun puisse contribuer sans casser le travail des autres.

Et dans mes projets universitaires en groupe, j'ai souvent pris le
r\^ole de structuration technique --- d\'efinir le cadre de travail
pour que chacun puisse avancer de mani\`ere autonome. \fg{}
\end{verbatimbox}

\vspace{8pt}

% ============================================================
\section{Questions techniques potentielles}
% ============================================================

\textcolor{warmgray}{\textit{L'interviewer est un manager CSPM, pas un quant. Les questions
techniques seront orient\'ees m\'etier, pas mod\'elisation.}}

\vspace{6pt}

\begin{graybox}{Q1 : \og Le LTV d'un client passe de 65\,\% \`a 85\,\%. Que faites-vous ? \fg{}}
\textbf{R\'eponse structur\'ee} :
\begin{enumerate}[nosep, leftmargin=16pt]
\item V\'erifier la cause : la NAV a baiss\'e ? Le client a tir\'e davantage sur sa ligne ?
\item Comparer au seuil contractuel (covenant) --- si on est au-dessus, margin call
\item Calculer le montant de l'appel de marge : combien faut-il pour revenir sous le seuil
\item Contacter le client et/ou escalader au manager
\item Documenter et informer le risk management
\end{enumerate}
\end{graybox}

\vspace{4pt}

\begin{graybox}{Q2 : \og Qu'est-ce qu'un haircut et pourquoi c'est important ? \fg{}}
\textbf{R\'eponse} : C'est la d\'ecote appliqu\'ee au collat\'eral pour absorber la volatilit\'e
en cas de liquidation. Un haircut de 20\,\% sur un portefeuille de 100\,M\texteuro{} signifie
qu'on ne compte que 80\,M\texteuro{} de garantie. C'est important parce que \c{c}a d\'etermine
directement la capacit\'e de financement du client et le LTV effectif.
\end{graybox}

\vspace{4pt}

\begin{graybox}{Q3 : \og Quelle diff\'erence entre financement ABC et financement classique ? \fg{}}
\textbf{R\'eponse} : Le financement classique repose sur la NAV globale du fonds.
L'ABC (Asset-Based Credit) regarde actif par actif --- chaque ligne a son propre
taux d'avance selon sa liquidit\'e et son risque. C'est plus granulaire, plus
conservateur, mais permet de financer des portefeuilles h\'et\'erog\`enes o\`u la NAV
globale serait trop volatile.
\end{graybox}

\vspace{4pt}

\begin{graybox}{Q4 : \og Comment l'IA peut aider dans le monitoring de portefeuilles ? \fg{}}
\textbf{R\'eponse} : Trois niveaux. D'abord, \textbf{automatisation des calculs r\'ep\'etitifs}
--- calcul de LTV quotidien, extraction de NAV, g\'en\'eration de reportings.
Ensuite, \textbf{d\'etection d'anomalies} --- rep\'erer les fonds dont le
comportement d\'evie de l'historique. Enfin, \textbf{acc\'el\'eration des analyses
ad hoc} --- quand un manager demande un nouveau reporting, l'IA permet de le
prototyper en heures plut\^ot qu'en semaines.
\end{graybox}

\vspace{4pt}

\begin{graybox}{Q5 : \og Expliquez IFRS\,9 en une minute \fg{}}
\textbf{R\'eponse} : IFRS\,9 est la norme comptable qui r\'egit le provisionnement
des cr\'eances. Avant (IAS\,39), on provisionnait quand la perte \'etait \og av\'er\'ee \fg{}.
Depuis IFRS\,9, on provisionne les pertes \og attendues \fg{} (ECL = PD $\times$ LGD $\times$ EAD).
Les cr\'eances sont class\'ees en 3 stages : Stage 1 (sain, perte \`a 12 mois),
Stage 2 (d\'egradation significative, perte \textit{lifetime}), Stage 3 (d\'efaut).
Pour CSPM, c'est pertinent parce qu'un client dont le LTV se d\'egrade peut passer
en Stage 2, ce qui augmente le co\^ut de provisionnement pour la banque.
\end{graybox}

\vspace{8pt}

% ============================================================
\section{Gestion du stress \& body language}
% ============================================================

\begin{goodbox}{Les fondamentaux}
\begin{itemize}[nosep, leftmargin=12pt]
\item \textbf{Contact visuel} --- alterner entre les deux interviewers (2 entretiens avec des managers op\'erationnels)
\item \textbf{Posture} --- droit, l\'eg\`erement pench\'e en avant (int\'er\^et), mains visibles sur la table
\item \textbf{Silences} --- c'est \textbf{OK} de prendre 3 secondes pour r\'efl\'echir avant de r\'epondre.
      Dire \og C'est une bonne question, laissez-moi structurer ma r\'eponse \fg{} est mieux qu'un \og euh\ldots \fg{}
\item \textbf{Si tu ne sais pas} --- \og Je ne connais pas ce produit en d\'etail, mais
      voici comment j'aborderais le probl\`eme\ldots \fg{}. Ne jamais inventer.
\item \textbf{Sourire} --- un entretien c'est aussi un test de compatibilit\'e humaine.
      L'\'equipe cherche quelqu'un avec qui ils veulent travailler 6 mois.
\end{itemize}
\end{goodbox}

\vspace{8pt}

% ============================================================
\section{Logistique \& checklist}
% ============================================================

\begin{graybox}{Avant l'entretien}
\begin{itemize}[nosep, leftmargin=12pt]
\item[$\square$] \textbf{Tenue} : costume ou blazer + chemise, pas de cravate (stage FO, pas interview MD).
      Pantin = bureau, pas trading floor, mais rester professionnel.
\item[$\square$] \textbf{Trajet} : v\'erifier l'acc\`es au b\^atiment BNP \`a Pantin (m\'etro ligne 5,
      Hoche ou Eglise de Pantin). Arriver 10--15 min en avance.
\item[$\square$] \textbf{CV imprim\'e} : 2--3 copies (un par interviewer + un pour toi).
\item[$\square$] \textbf{Bloc-notes} : pour noter les noms des interviewers et leurs questions.
\item[$\square$] \textbf{Recherche LinkedIn} : identifier les managers qui t'interviewent si possible.
      Conna\^itre leur parcours montre de l'int\'er\^et.
\end{itemize}
\end{graybox}

\begin{graybox}{Apr\`es l'entretien}
\begin{itemize}[nosep, leftmargin=12pt]
\item[$\square$] \textbf{Email de remerciement} dans les 24h --- court (5 lignes), personnalis\'e
      avec un \'el\'ement sp\'ecifique de la conversation.
\item[$\square$] \textbf{D\'ebriefing personnel} : noter les questions pos\'ees, tes r\'eponses,
      ce qui a bien march\'e et ce qui peut \^etre am\'elior\'e (utile si deuxi\`eme tour).
\end{itemize}
\end{graybox}

\vspace{8pt}

% ============================================================
\section{Questions \`a poser en fin d'entretien}
% ============================================================

\textcolor{warmgray}{\textit{Pr\'epare 4--5 questions, poses-en 2--3 selon le temps. \'Evite
les questions dont la r\'eponse est dans l'annonce.}}

\begin{enumerate}[leftmargin=16pt]
\item \og Quels types de fonds repr\'esentent la majorit\'e de l'activit\'e CSPM
      aujourd'hui --- plut\^ot hedge funds, fonds de dette, ou mixte ? \fg{}
\item \og Comment se passe concr\`etement la coordination avec l'\'equipe de
      recherche \`a New York ? \fg{}
\item \og Quel est l'outil principal utilis\'e pour le suivi de LTV au quotidien
      --- d\'eveloppement interne, Bloomberg, ou autre ? \fg{}
\item \og Y a-t-il des projets d'automatisation en cours dans l'\'equipe o\`u un
      stagiaire pourrait apporter de la valeur ajout\'ee ? \fg{}
\item \og Quel est le plus grand d\'efi op\'erationnel de l'\'equipe en ce moment ? \fg{}
\end{enumerate}

\vspace{8pt}

% ============================================================
\section{Pi\`eges \`a \'eviter --- R\'ecapitulatif}
% ============================================================

\begin{dangerbox}{Les 10 erreurs qui tuent un entretien}
\begin{enumerate}[nosep, leftmargin=16pt]
\item \textbf{R\'eciter} un pitch appris par c\oe{}ur --- adapt\'e \`a la dynamique
\item \textbf{Parler trop longtemps} --- r\`egle des 90 secondes par r\'eponse
\item \textbf{Dire} \og j'ai chang\'e de projet \fg{} --- dis \og j'ai approfondi \fg{}
\item \textbf{Sur-vendre la technique} (TabNet, Mahalanobis\ldots) \`a un manager Front Office
\item \textbf{Minimiser le m\'emoire M1} --- c'est lui le match direct
\item \textbf{Dire} \og c'est un projet acad\'emique \fg{} --- dis \og j'ai construit / impl\'ement\'e \fg{}
\item \textbf{Mentir sur l'IA} --- l'honn\^etet\'e + le framing \og pilotage \fg{} est bien plus fort
\item \textbf{Critiquer} une ancienne \'equipe/projet/prof --- toujours positif
\item \textbf{Ne poser aucune question} --- montre un manque d'int\'er\^et
\item \textbf{Parler salaire en premier} --- c'est un stage, la gratification est fix\'ee par la loi
\end{enumerate}
\end{dangerbox}

\vspace{10pt}

% %%%%%%%%%%%%%%%%%%%%%%%%%%%%%%%%%%%%%%%%%%%%%%%%%%%%%%%%%%%%
%
%         PARTIE III --- ENGLISH SWITCH
%
% %%%%%%%%%%%%%%%%%%%%%%%%%%%%%%%%%%%%%%%%%%%%%%%%%%%%%%%%%%%%

\newpage
\part*{Partie III --- English switch}
\addcontentsline{toc}{part}{Partie III --- English switch}

% ============================================================
\section{G\'erer le passage en anglais}
% ============================================================

\begin{warnbox}{Pourquoi c'est tr\`es probable}
\begin{itemize}[nosep, leftmargin=12pt]
\item L'annonce exige \textbf{\og anglais courant \fg{}}
\item L'\'equipe CSPM coordonne avec \textbf{New York et Londres} au quotidien
\item Le manager CSPM travaille avec l'\'equipe de recherche sur les fonds sp\'ecialis\'es \`a NY
\item En CIB, le switch en anglais en milieu d'entretien est un \textbf{classique}
      pour tester le niveau r\'eel --- plus fiable qu'un score TOEIC
\end{itemize}
\end{warnbox}

\begin{goodbox}{Comment r\'eagir au switch}
\begin{itemize}[nosep, leftmargin=12pt]
\item \textbf{Ne pas paniquer.} Prends une seconde, souris, et r\'eponds naturellement.
\item \textbf{Ne t'excuse pas} pour ton anglais --- ne dis pas \og Sorry, my English is not perfect \fg{}.
      Parle, c'est tout.
\item \textbf{Accepte un vocabulaire plus simple.} En anglais, tes phrases seront plus courtes
      et moins nuanc\'ees qu'en fran\c{c}ais --- c'est normal et \c{c}a suffit.
\item \textbf{Utilise le jargon financier.} LTV, NAV, margin call, haircut, covenant, ECL, staging
      --- ces mots sont les m\^emes dans les deux langues. C'est ton filet de s\'ecurit\'e.
\item \textbf{Si tu bloques sur un mot} : d\'ecris le concept. \og The ratio that measures how much
      the loan is covered by the collateral \fg{} est parfaitement valide si tu oublies \og LTV \fg{}.
\end{itemize}
\end{goodbox}

\vspace{8pt}

% ============================================================
\section{Pitch --- English version ($\sim$90 seconds)}
% ============================================================

\begin{pitchbox}
``My Master's thesis at Paris-Saclay focuses on \textbf{fund monitoring and
collateral management}: I built a decision-support system for hedge funds,
with automated NAV and LTV calculations, margin call tracking, collateral
valuation with haircuts, and liquidity stress tests on portfolios.

Alongside this, I developed a full \textbf{IFRS\,9 credit risk cockpit} ---
ECL computation, probability of default models, macroeconomic stress testing,
RAROC and capital allocation --- which gave me an \textbf{end-to-end view}
of the risk chain, from individual exposures all the way up to CRO-level
reporting.

Both projects converge towards the CSPM role: the first one gave me the
\textbf{Front Office toolkit} --- LTV, collateral, margin calls --- and the
second one gave me the \textbf{analytical depth on the Risk side}. That's
exactly the intersection of this position, at the boundary between Front
Office and the Risk department.''
\end{pitchbox}

\vspace{8pt}

% ============================================================
\section{Key questions --- English answers}
% ============================================================

\subsection{``Tell me about yourself''}

\begin{verbatimbox}
``I'm a Master's student at Paris-Saclay, specialising in quantitative
finance and economics. My work is oriented towards risk analysis and
financing.

\textit{[Insert the 90-second pitch above]}

I'm looking for an internship where I can apply these skills in an
operational environment, with direct client interaction. The CSPM role
is exactly that.''
\end{verbatimbox}

\subsection{``Why BNP Paribas? Why CSPM?''}

\begin{verbatimbox}
``Three reasons. First, the role sits at the exact intersection of my
two projects --- Front Office with a strong risk dimension. Second,
the CSPM team covers a wide range of financing products --- hedge funds,
ABC, Capdoc, fund of funds --- which gives me broad exposure to the
fund financing business. And third, BNP's Equity and Prime Services
division is the fastest-growing segment in CIB --- over 10 percent
revenue growth in 2025, and prime brokerage balances have quadrupled
in six years. It's an exciting time to join this business.''
\end{verbatimbox}

\subsection{``What is LTV and why does it matter?''}

\begin{verbatimbox}
``LTV stands for Loan-to-Value. It's the ratio of the outstanding loan
to the value of the collateral. For example, if a client has a 70-million
euro loan secured by a portfolio worth 100 million, the LTV is 70 percent.

It matters because it's the key risk indicator for fund financing. When
the LTV goes up --- either because the collateral value drops or the client
draws more on the facility --- the bank's exposure increases relative to the
guarantee. If it breaches the contractual covenant, we need to issue a
margin call.''
\end{verbatimbox}

\subsection{``Walk me through a margin call scenario''}

\begin{verbatimbox}
``Let's say a client's LTV goes from 65 to 85 percent, and the covenant
is set at 75 percent. First, I'd check the cause --- has the NAV of the
fund dropped, or did the client draw more on the facility?

Then I'd calculate how much additional collateral is needed to bring the
LTV back below 75 percent. I'd prepare the analysis for the manager, and
we'd contact the client to request either additional collateral or a
partial repayment of the loan. At the same time, I'd flag it to the
risk department and document everything.''
\end{verbatimbox}

\subsection{``Tell me about your IFRS\,9 project''}

\begin{verbatimbox}
``I built a complete IFRS\,9 credit risk cockpit in Python. It covers the
full pipeline: PD modelling with three model families --- logistic regression
with scorecards, gradient boosting, and deep learning --- LGD estimation
with cyclical adjustments, and a multi-factor SICR staging engine.

The cockpit also computes RAROC, capital allocation, reverse stress tests,
and generates a 7-tab dashboard for CRO reporting. The whole portfolio
recalculates in about one second.

What's directly relevant for CSPM is that I understand how the risk
assumed by the bank on each exposure relates to the collateral held.
When the team discusses with the risk department, I can speak their
language.''
\end{verbatimbox}

\subsection{``How did you build a 15,000-line project as an M1 student?''}

\begin{verbatimbox}
``I used AI as an implementation accelerator --- not just basic prompts,
but a structured methodology. I designed the architecture upfront, wrote
detailed functional specifications, then steered the AI for implementation,
and most importantly, I reviewed, validated and corrected every module.

My value added is not writing Python code. It's knowing that an ECL should
be computed using hazard rates rather than a Bernoulli approximation, that
LGD needs to be cyclical, that RAROC must include a cost-income ratio and
a tax rate. The AI implements fast, but it doesn't know whether the output
is financially correct. I do.

For a team like CSPM, it means that if you need to automate a new LTV
tracker or build a reporting tool, I can deliver in days rather than weeks.
It's a real productivity lever.''
\end{verbatimbox}

\subsection{``What do you know about BNP Paribas?''}

\begin{verbatimbox}
``BNP Paribas is the largest bank in the eurozone, with a record net income
of 12.2 billion euros in 2025. On the CIB side, Equity and Prime Services
is the fastest-growing business line --- over 10 percent growth last year ---
driven by the Deutsche Bank prime brokerage acquisition in 2019 and the
capture of clients after Credit Suisse exited in 2023.

What interests me specifically is that BNP is one of only four banks
globally that offer the full suite of investor services --- prime brokerage,
clearing, listed derivatives, and securities financing. And with over
800 AI applications in production and a partnership with Mistral AI,
the bank is also investing heavily in technology.''
\end{verbatimbox}

\subsection{``What are your strengths and weaknesses?''}

\begin{verbatimbox}
``I'm rigorous --- my IFRS\,9 cockpit has 600 automated tests and went
through three mathematical audits. I'm also autonomous --- both projects
were personal initiatives, not coursework.

My area for improvement is that I sometimes tend to over-optimise before
delivering. I'm learning to prioritise pragmatism --- ship a working
version first, then iterate.''
\end{verbatimbox}

\subsection{``Where do you see yourself in five years?''}

\begin{verbatimbox}
``In the short term, I want to build a solid foundation in fund financing
--- understand the products, the client relationships, the day-to-day
monitoring. If the internship goes well, the Graduate Programme or a VIE
would let me see other desks or other geographies.

In the medium term, I see myself growing into a senior analyst role in
fund financing, at the intersection of structuring and risk management ---
exactly where CSPM sits today.''
\end{verbatimbox}

\vspace{8pt}

% ============================================================
\section{English vocabulary cheat sheet}
% ============================================================

\textcolor{warmgray}{\textit{Termes que tu utilises en fran\c{c}ais et dont tu dois conna\^itre
l'\'equivalent anglais naturel.}}

\begin{center}
\small
\begin{tabular}{@{}ll|ll@{}}
\toprule
\textbf{Fran\c{c}ais} & \textbf{English} & \textbf{Fran\c{c}ais} & \textbf{English} \\
\midrule
Collat\'eral & Collateral & Appel de marge & Margin call \\
D\'ecote & Haircut & Covenant & Covenant \\
Encours & Outstanding / Exposure & Ligne de cr\'edit & Credit facility \\
Taux d'avance & Advance rate & Seuil & Threshold / Trigger \\
Fonds de couverture & Hedge fund & Fonds de fonds & Fund of funds \\
Engagement non appel\'e & Unfunded commitment & Levier & Leverage \\
Tirage & Drawdown & \'Ech\'eancier & Schedule / Amortisation \\
Provisionnement & Provisioning & Perte attendue & Expected loss \\
D\'egradation & Deterioration & Portefeuille & Portfolio \\
Rendement & Return / Yield & Actif pond\'er\'e & Risk-weighted asset \\
B\^ale / CRR & Basel / CRR & Ratio de solvabilit\'e & Capital adequacy ratio \\
Structuration & Structuring & R\`eglement-livraison & Settlement \\
Marge commerciale & Commercial spread & Co\^ut du risque & Cost of risk \\
Notation & Rating & D\'efaut & Default \\
Flux de tr\'esorerie & Cash flow & Valeur liquidative & Net Asset Value (NAV) \\
\bottomrule
\end{tabular}
\end{center}

\vspace{6pt}

\begin{warnbox}{Faux amis \& pi\`eges linguistiques}
\begin{itemize}[nosep, leftmargin=12pt]
\item \textbf{``Actually''} $\neq$ actuellement. Actuellement = \textit{currently}.
      Actually = \textit{en fait}.
\item \textbf{``Eventually''} $\neq$ \'eventuellement. Eventually = \textit{finalement / \`a terme}.
      \'Eventuellement = \textit{possibly}.
\item \textbf{``Delay''} $\neq$ d\'elai. Delay = retard. D\'elai = \textit{deadline / timeframe}.
\item \textbf{``Resume''} $\neq$ r\'esum\'e. Resume = reprendre. R\'esum\'e (en anglais) = CV.
      R\'esum\'e (en fran\c{c}ais) = \textit{summary}.
\item \textbf{``Sensible''} $\neq$ sensible. Sensible (EN) = \textit{raisonnable}.
      Sensible (FR) = \textit{sensitive}.
\item \textbf{``Comprehensive''} $\neq$ compr\'ehensif. Comprehensive = \textit{exhaustif / complet}.
\end{itemize}
\end{warnbox}

\vspace{6pt}

\begin{goodbox}{Phrases de transition utiles}
\begin{itemize}[nosep, leftmargin=12pt]
\item \textit{``That's a great question. Let me walk you through...''}
      --- pour gagner 2 secondes de r\'eflexion
\item \textit{``To put it simply...''} --- pour simplifier un concept technique
\item \textit{``What I find particularly relevant for CSPM is that...''}
      --- pour toujours relier au poste
\item \textit{``In practical terms, this means...''}
      --- pour passer de la th\'eorie \`a l'op\'erationnel
\item \textit{``I'm not entirely familiar with that specific product, but
      here's how I'd approach it...''} --- si tu ne sais pas
\item \textit{``From a risk perspective...''} --- pour montrer ta double casquette
\item \textit{``The way I see it...''} --- pour donner ton point de vue
\end{itemize}
\end{goodbox}

\vspace{10pt}

% ============================================================
\section{Aide-m\'emoire express --- \`A relire 30 min avant}
% ============================================================

\begin{goodbox}{Les 7 messages \`a faire passer}
\begin{enumerate}[nosep, leftmargin=16pt]
\item Mon m\'emoire M1 porte \textbf{exactement} sur le monitoring de fonds et le collat\'eral
\item Le cockpit IFRS\,9 me donne la \textbf{profondeur risque} que les autres n'ont pas
\item Je sais utiliser l'IA comme \textbf{acc\'el\'erateur}, pas comme b\'equille
\item Je suis \textbf{op\'erationnel rapidement} : LTV, margin calls, Python, Excel
\item Je comprends le \textbf{langage des risques} (ECL, staging, RWA)
\item Je suis \textbf{rigoureux} (600 tests, 3 audits) et \textbf{autonome}
\item Je veux apprendre le \textbf{m\'etier du financement} sur le terrain
\end{enumerate}
\end{goodbox}

\vspace{12pt}

\begin{center}
\textcolor{warmgray}{\rule{0.4\textwidth}{0.5pt}}\\[6pt]
\textcolor{warmgray}{\small\textit{Structure : pitch 90s + modules \`a la demande.
Partie II = filet de s\'ecurit\'e, pas un script.\\
Laisser l'interviewer guider. \^Etre naturel. Montrer l'envie.}}
\end{center}

\end{document}
